\section{Planets and Kingdoms}
\label{sec:cses-planets-kingdoms}

\problemstatement{Given a directed graph of $n$ planets and $m$
teleporters, partition the planets into kingdoms so that two planets are in
the same kingdom iff each can reach the other. Output the number of
kingdoms and each planet's kingdom id.}

\begin{patternbox}
``Mutually reachable groups'' are the \textbf{strongly connected components}
(SCCs) of a directed graph. \textbf{Tarjan's algorithm} finds all SCCs in a
single DFS using discovery times and a low-link value that detects the root
of each component.
\end{patternbox}

\subsection*{Tarjan's low-link}

DFS assigns each node an increasing discovery index. A node's
$\textit{low}$ is the smallest index reachable from it using tree edges
plus at most one back edge to a node still on the DFS stack. When a node's
$\textit{low}$ equals its own index, it is the root of an SCC: pop the DFS
stack down to it, and those nodes form one component. Nodes not yet assigned
sit on an explicit stack so cross-component edges are excluded.

\begin{ideabox}[low == index Marks a Component Root]
The low-link captures ``the earliest ancestor I can loop back to.'' If a
node cannot reach anything older than itself, it closes a strongly connected
region, so everything above it on the stack is exactly that SCC.
\end{ideabox}

\subsection*{Algorithm}

\begin{enumerate}
  \item DFS every node, maintaining discovery indices, low-links, and an
        on-stack marker.
  \item When $\textit{low}[u]=\textit{index}[u]$, pop the stack to $u$,
        labelling those nodes with a new component id.
  \item Output the component count and each node's id.
\end{enumerate}

\subsection*{Dry run}

\dryrun{$1\!\to\!2$, $2\!\to\!1$, $2\!\to\!3$.}

\begin{itemize}
  \item $\{1,2\}$ are mutually reachable -- one kingdom; $\{3\}$ alone.
  \item Two kingdoms.
\end{itemize}

\subsection*{C++ implementation}

\begin{lstlisting}
#include <bits/stdc++.h>
using namespace std;

int n, m, timer_ = 0, sccCount = 0;
vector<vector<int>> adj;
vector<int> disc, low, comp;
vector<bool> onStack;
stack<int> stk;

void tarjan(int u) {
    disc[u] = low[u] = ++timer_;
    stk.push(u); onStack[u] = true;
    for (int v : adj[u]) {
        if (!disc[v]) { tarjan(v); low[u] = min(low[u], low[v]); }
        else if (onStack[v]) low[u] = min(low[u], disc[v]);
    }
    if (low[u] == disc[u]) {                      // u is an SCC root
        sccCount++;
        while (true) {
            int x = stk.top(); stk.pop(); onStack[x] = false;
            comp[x] = sccCount;
            if (x == u) break;
        }
    }
}

int main() {
    cin >> n >> m;
    adj.assign(n + 1, {});
    for (int i = 0; i < m; i++) { int a, b; cin >> a >> b; adj[a].push_back(b); }
    disc.assign(n + 1, 0); low.assign(n + 1, 0);
    comp.assign(n + 1, 0); onStack.assign(n + 1, false);

    for (int i = 1; i <= n; i++) if (!disc[i]) tarjan(i);

    cout << sccCount << "\n";
    for (int i = 1; i <= n; i++) cout << comp[i] << " ";
    cout << "\n";
    return 0;
}
\end{lstlisting}

\begin{complexitybox}
\textbf{Time Complexity.} $O(n+m)$ -- a single DFS. \textbf{Space
Complexity.} $O(n+m)$; recursion depth $O(n)$ (use an iterative DFS for very
large $n$).
\end{complexitybox}

\begin{takeawaybox}
Mutually-reachable partitions are SCCs, and Tarjan finds them in one DFS via
the low-link ``earliest reachable ancestor'' test. SCCs underlie strong-
connectivity checks, 2-SAT, and condensation DP -- the next three problems.
\end{takeawaybox}
